%% kapa: exam=1 round=27 year=2016 period=1 form=A subject=real_estate subject_ko=부동산학원론 seq=1 number=81
%% source: Q-Net 감정평가사 기출문제 (cst00309) artlSeq=5169402 file=extracted/1008/5169402/제27회 감정평가사 1차 문제/제27회 감정평가사 1교시문제(A형).pdf
%% conversion: kordoc 4.24.51 -> tools/gen1.py (auto); review=proofread-sonnet; flags=none
\begin{question}{81}
부동산의 개념에 관한 설명으로 옳지 않은 것은?
\begin{choices}
  \item 법률적 개념에서 협의의 부동산은 민법 제99조 제1항에서의 ‘토지 및 그 정착물’을 말한다.
  \item 부동산의 경우에는 등기로써 공시의 효과를 가지지만 동산은 점유로써 공시의 효과를 가진다.
  \item 좁은 의미의 부동산과 준부동산을 합쳐 광의의 부동산이라 하며, 자본, 자산 등과 함께 기술적 측면에서의 부동산으로 구분된다.
  \item 준부동산은 물권변동을 등기나 등록수단으로 공시하는 동산을 포함한다.
  \item 입목에 관한 법령에 의해 소유권보존등기된 입목, 공장 및 광업재단 저당법령에 의하여 저당권의 목적물이 되고 있는 공장재단은 부동산에 준하여 취급된다.
\end{choices}
\end{question}
\answer{3}
